% !TeX root = Javier_TARAZONA_CV_fr.tex
% !TeX program = pdflatex
%-------------------------
% Resume in Latex
% Author : Javier Tarazona
% Offre : Renault Group - Ingenieur IA & Data Industrielle - Intelligence Multimodale pour l'Analyse Video en Manufacturing (Workday)
%------------------------

\documentclass[a4paper,10pt]{article}
\DeclareUnicodeCharacter{00A0}{~}
\DeclareUnicodeCharacter{2013}{\textendash}
\DeclareUnicodeCharacter{2014}{\textemdash}
\DeclareUnicodeCharacter{201C}{``}
\DeclareUnicodeCharacter{201D}{''}
\DeclareUnicodeCharacter{2022}{\textbullet}
\DeclareUnicodeCharacter{25E6}{\textopenbullet}
\usepackage[T1]{fontenc}      % Codificación de salida (acentos y símbolos)
\usepackage[utf8]{inputenc}   % Lee texto en UTF-8
\usepackage[empty]{fullpage}
\usepackage[export]{adjustbox}
\usepackage[pdftex]{hyperref}
\usepackage[usenames,dvipsnames]{color}
\usepackage{array}
\usepackage{comment}
\usepackage{enumitem}         % Control de listas
\usepackage{fancyhdr}
\usepackage{graphicx}
\usepackage{latexsym}
\usepackage{lmodern}
\usepackage{marvosym}
\usepackage{tabularx}
\usepackage{titlesec}
\usepackage{verbatim}
\usepackage{xcolor}

% --- Viñetas limpias y seguras ---
\setlist[itemize]{label=\textbullet, labelsep=0.5em, leftmargin=*}

% ===== CORRECTIFS ATS =====
\input{glyphtounicode}        % chaque glyphe -> vrai texte Unicode (ligatures decomposees)
\pdfgentounicode=1
\usepackage[none]{hyphenat}   % point 5 : plus aucun mot-cle coupe
\sloppy                       %           evite les debordements de marge
\usepackage{textcomp}         % point 6 : apostrophe ASCII au lieu de U+2019
\catcode`\'=\active
\def'{\textquotesingle}
\hypersetup{                  % point 8 : metadonnees indexables
  pdftitle={Javier Andres TARAZONA JIMENEZ - CV - Stage de fin d'études - Ingénieur IA \& Data Industrielle - Intelligence Multimodale pour l'Analyse Vidéo en Manufacturing},
  pdfauthor={Javier Andres TARAZONA JIMENEZ},
  pdfsubject={Stage de fin d'études Bac+5 - Ingénieur IA \& Data Industrielle - IA multimodale, analyse vidéo, IA générative, LLM},
  pdfkeywords={Ingénieur IA, Data Industrielle, Intelligence Multimodale, Analyse Vidéo, Manufacturing,
               IA générative, Large Language Models (LLM), Vision Language Models (VLM), Vision par ordinateur,
               Computer Vision, Python, PyTorch, OpenCV, Deep Learning, Data Science, Vision Transformers,
               Bac+5, Diplôme d'Ingénieur, Grande École, stage de fin d'études},
  pdfborder={0 0 0}           % pas de cadre autour des liens (le soulignement bleu suffit)
}
% ==========================

\pagestyle{empty}

% Adjust margins
\addtolength{\oddsidemargin}{-0.45in}
\addtolength{\evensidemargin}{-0.45in}
\addtolength{\textwidth}{1.15in}
\addtolength{\topmargin}{-.62in}
\addtolength{\textheight}{1.20in}

\urlstyle{same}

\raggedbottom
\raggedright
\setlength{\tabcolsep}{0in}
\renewcommand{\arraystretch}{0.98}

%\definecolor{mydarkgreen}{RGB}{0, 140, 123}
\definecolor{mydarkgreen}{RGB}{0, 0, 0}

% Sections formatting
\titleformat{\section}{
  \vspace{-6pt}\scshape\raggedright\large
}{}{0em}{}[\color{mydarkgreen}\titlerule \vspace{-8pt}]

%-------------------------
% Custom commands

% Indented paragraph helper: \indentbox[ancho]{contenido}
% Uso: \indentbox{Texto...}  o  \indentbox[1.5em]{Texto...}
% [t] + \raggedright + \strut : lignes alignees en haut, sans espaces etires ni chevauchement.
\newcommand{\indentbox}[2][1em]{\hspace*{#1}\parbox[t]{\dimexpr\linewidth-#1\relax}{\raggedright #2\strut}}

% Photo et QR code desactives : dans Workday une photo casse l'analyse du texte voisin.


%-------------------------------------------
%%%%%%  CV STARTS HERE  %%%%%%%%%%%%%%%%%%%%%%%%%%%%


\begin{document}

%----------HEADING-----------------
% Nom seul sur la premiere ligne, puis coordonnees en texte brut dans le corps (pas d'en-tete, pas de minipage).

\noindent{\Large\bfseries\textcolor{mydarkgreen}{Javier Andres TARAZONA JIMENEZ}}\par
\vspace{3pt}
\noindent{\footnotesize javitar06@gmail.com \textbullet{} +33 7 46 36 97 75 \textbullet{} Massy, France}\par
\vspace{1pt}
\noindent{\footnotesize\href{https://www.linkedin.com/in/javier-andres-tarazona-jimenez-84b489222/}{\textcolor{blue}{\underline{linkedin.com/in/javier-andres-tarazona-jimenez-84b489222}}} \textbullet{} \href{https://github.com/JavierTarazona06}{\textcolor{blue}{\underline{github.com/JavierTarazona06}}}}\par

\vspace{0.12cm}

{\centering
  {\large\bfseries Stage de fin d'études - Ingénieur IA \& Data Industrielle}\\[1pt]
  {\normalsize\bfseries Intelligence Multimodale pour l'Analyse Vidéo en Manufacturing}\\[2pt]
  {\footnotesize Élève en dernière année d'école d'ingénieur (Grande École, Bac+5) \textbullet{} Disponible à partir d'avril 2027 pour 6 mois}\\[3pt]
  {\footnotesize\textit{Vision par ordinateur (Vision Transformers, analyse vidéo) et LLM : les briques de l'IA multimodale et des Vision Language Models (VLM)}}\par}

\vspace{0.05cm}

%-----------FORMATION-----------------
\section{\textcolor{mydarkgreen}{\textbf{FORMATION}}}
{\raggedright
\textbf{ENSTA Paris (Institut Polytechnique de Paris)} \textbullet{} Palaiseau, France \textbullet{} \textit{\footnotesize 2025/09 - 2027/09}\\
{\footnotesize\textit{Diplôme d'Ingénieur (Bac+5) - Master of Science in Engineering, parcours Intelligence Artificielle}}\\[3pt]
{\footnotesize Cours : Machine Learning, Reconnaissance d'images, Bases de données, Ingénierie des systèmes, Économie industrielle et innovation}\\[4pt]

\textbf{Universidad Nacional de Colombia (UNAL)} \textbullet{} Bogotá, Colombie \textbullet{} \textit{\footnotesize 2021/01 - 2025/07}\\
{\footnotesize\textit{Cursus d'ingénieur - Ingénierie des Systèmes et Informatique}}\\[3pt]
{\footnotesize 12e université d'Amérique latine (QS Latin America 2026). Cours : Probabilités et statistique, Modèles stochastiques, Optimisation}\\[4pt]

\textbf{University of Saskatchewan} \textbullet{} Saskatoon, Canada \textbullet{} \textit{\footnotesize 2024/09 - 2024/12}\\
{\footnotesize\textit{Échange académique - Informatique : Ingénierie logicielle, Algorithmes, Systèmes d'exploitation}}\par}

\vspace{0.05cm}
%-----------EXPERIENCE-----------------
\section{\textcolor{mydarkgreen}{\textbf{EXPÉRIENCE PROFESSIONNELLE}}}
{\raggedright
\textbf{ENSTA Paris - Laboratoire U2IS} \textbullet{} Palaiseau, France \textbullet{} \textit{\footnotesize 2026/05 - 2026/08}\\
{\footnotesize\textit{Stagiaire de recherche en Intelligence Artificielle et Vision 3D}}\\[3pt]
\indentbox{\footnotesize \textbf{- Préparation de données :} pipeline Python/Blender générant des datasets d'entraînement 3D multi-vues (caméras, métadonnées).}\\
\indentbox{\footnotesize \textbf{- Deep Learning :} pré-entraînement par prédiction de vues avec \textbf{PyTorch} et \textbf{Vision Transformers} pour apprendre la structure 3D.}\\
\indentbox{\footnotesize \textbf{- Vision 3D auto-supervisée :} synthèse de nouvelles vues sur formes géométriques de complexité croissante (curriculum learning).}\\[4pt]

\textbf{APEDYS91} \textbullet{} Paris, France \textbullet{} \textit{\footnotesize 2025/10 - 2026/02}\\
{\footnotesize\textit{Ingénieur Machine Learning / NLP}}\\[3pt]
\indentbox{\footnotesize \textbf{- IA générative et LLM :} correction de texte en français, pipeline hybride LanguageTool et \textbf{Large Language Models (LLM)} locaux.}\\
\indentbox{\footnotesize \textbf{- Fiabilité et intégration :} guardrails NER (spaCy) et contrôle des entités/nombres anti-hallucination ; modèle adapté à la RAM/VRAM.}\\[4pt]

\textbf{Engin A.I} \textbullet{} Bogotá, Colombie (à distance, États-Unis) \textbullet{} \textit{\footnotesize 2025/02 - 2025/06}\\
{\footnotesize\textit{Ingénieur Logiciel et Vision par Ordinateur}}\\[3pt]
\indentbox{\footnotesize \textbf{- Vision par ordinateur} (OpenCV, NumPy) \textbf{:} \textbf{+50\% de vitesse et +60\% de précision} du post-traitement de détection des voies et marquages routiers ; clustering non supervisé pour la robustesse.}\\
\indentbox{\footnotesize \textbf{- Analyse de données :} trajectoires de roue (WheelPath) corrélées aux zones de détérioration (+10\% de satisfaction utilisateur).}\\
\indentbox{\footnotesize \textbf{- Architecture logicielle Python :} refactorisation modulaire de 2 bibliothèques, \textbf{+70\% de score qualité} (maintenabilité, tests).}\\[4pt]

\textbf{Engin A.I} \textbullet{} Bogotá, Colombie (à distance, États-Unis) \textbullet{} \textit{\footnotesize 2024/01 - 2024/07}\\
{\footnotesize\textit{Ingénieur Logiciel et Vision par Ordinateur}}\par}


\vspace{0.05cm}
%-----------PROJETS-----------------
\section{\textcolor{mydarkgreen}{\textbf{PROJETS}}}
{\raggedright
\textbf{"Tameo" - Vision embarquée pour bateau autonome} \textbullet{} \textit{\footnotesize 2025/09 - 2027/06}\\
{\footnotesize\textit{Équipe de 5, membre puis chef de projet depuis 2026/09 - Monaco Energy Boat Challenge, AI Class - Python, PyTorch, YOLO, caméra ZED}}\\
\indentbox{\footnotesize Pipeline de vision (détection et segmentation) pour la décision de mouvement ; préparation de datasets spécialisés, data augmentation, fine-tuning et \textbf{comparaison de modèles} : mAP, F1 et précision/rappel portés d'environ 80\% à 90\%.}\\[3pt]

\textbf{Slow - Application d'analyse vidéo du trafic} \textbullet{} \href{https://github.com/JavierTarazona06/slow}{\textcolor{blue}{\underline{github.com/JavierTarazona06/slow}}}\\
\indentbox{\footnotesize \textbf{Analyse vidéo} du trafic (équipe de 5, 6 mois) : suivi de véhicules et estimation de leur vitesse avec OpenCV, historique stocké dans MySQL.}\\[3pt]

\textbf{Sharp Sight - Pipeline de collecte et de traitement des données} \textbullet{} \href{https://github.com/JavierTarazona06/SharpSight}{\textcolor{blue}{\underline{github.com/JavierTarazona06/SharpSight}}}\\
\indentbox{\footnotesize Collecte automatisée d'offres e-commerce (Selenium), nettoyage, structuration et catégorisation des données (pandas), API REST FastAPI.}\par}


\vspace{0.05cm}
%-----------COMPETENCES TECHNIQUES-----------------
\section{\textcolor{mydarkgreen}{\textbf{COMPÉTENCES TECHNIQUES}}}
{\footnotesize
\setlength{\parindent}{0pt}
\hangindent=1.4em\hangafter=1 \textbf{IA et Data Science :} Machine Learning, Deep Learning, Vision par ordinateur (Computer Vision), analyse vidéo, détection d'objets, segmentation, suivi d'objets, IA générative, Large Language Models (LLM), NLP, fine-tuning, data augmentation, évaluation de modèles\par
\hangindent=1.4em\hangafter=1 \textbf{Écosystème IA/Data :} Python, PyTorch, scikit-learn, OpenCV, YOLO, Vision Transformers, NumPy, Pandas, Llama, Mistral, Gemma\par
\hangindent=1.4em\hangafter=1 \textbf{Données et industrialisation :} SQL (MySQL, PostgreSQL), Docker, FastAPI, API REST, Google Cloud (GCP), CI/CD, Git/GitHub\par
\hangindent=1.4em\hangafter=1 \textbf{Programmation et méthodes :} C++/C, Java, JavaScript, architecture logicielle, Agile (SCRUM)\par
}

\vspace{0.05cm}
%-----------LANGUES-----------------
\section{\textcolor{mydarkgreen}{\textbf{LANGUES}}}
{\footnotesize
\noindent \textbf{Anglais :} courant (C1 - IELTS 2024) \textbullet{} \textbf{Français :} avancé (B2 - DELF 2024) \textbullet{} \textbf{Espagnol :} langue maternelle \textbullet{} \textbf{Portugais :} notions
}

\vspace{0.05cm}
%--------DISTINCTIONS------------
\section{\textcolor{mydarkgreen}{\textbf{DISTINCTIONS}}}
{\footnotesize
\noindent \textbf{Bourse Eiffel Excellence} \textbullet{} bourse d'excellence du gouvernement français pour le Master à ENSTA Paris\\[1pt]
\noindent \textbf{Mention d'Honneur - ICPC Colombie 2023} \textbullet{} XXXVII Marathon National de Programmation ACIS REDIS\\[1pt]
\noindent \textbf{Excellence académique} \textbullet{} UNAL, exonération des frais de scolarité
}


%-------------------------------------------
\end{document}
